\chapter{Conclusiones}\label{ch:conclusiones}

El hallazgo central de esta tesis, en cumplimiento del objetivo general de caracterizar la evolución estructural del viento a 10 metros y la precipitación en los ciclones extratropicales del Atlántico Sudoccidental, es que la intensidad dinámica del vórtice no predice de forma lineal ni la magnitud ni la ubicación de sus impactos en superficie. La organización espacial del viento y de la precipitación evoluciona de manera cualitativamente distinta entre la Población Global y el subconjunto de mayor vorticidad (p90), y en los sistemas más intensos ambos campos terminan por ocupar cuadrantes geométricamente opuestos del dominio lagrangiano durante la madurez. El viento se concentra al noroeste, gobernado por la cinta transportadora fría (CCB) y la fractura frontal, mientras la precipitación queda confinada a una banda periférica enroscada al sureste y sur, orbitando el margen de la intrusión seca. Este desacople, espacial y, como se muestra más adelante, también temporal, no es una yuxtaposición casual de dos fenómenos independientes, sino una firma estructural robusta, regionalmente modulada y estadísticamente verificable. Los apartados siguientes reconstruyen cómo se llegó a esta conclusión, siguiendo el hilo argumental del ciclo de vida, del método y de la región, más que un recuento secuencial de objetivos.

\section{De la organización difusa a la reorganización ocluida}

La base de datos lagrangiana (Objetivo Específico 1) se construyó filtrando, sobre el catálogo de trayectorias y la clasificación objetiva de fases ya generada mediante \textit{CycloPhaser} \citep{coutodesouza2024new, deSouza2025CycloPhaser}, únicamente los sistemas con ciclo de vida completo, y estratificando la muestra resultante por intensidad dinámica ($\zeta_{850}^{\min}$) en Población Global y subconjunto p90. Esta construcción permitió aislar con precisión la condición estructural de cada sistema en cada fase, verificando que el 83\% de los ciclones p90 alcanza su mínimo de vorticidad durante la madurez y no antes o después. Esta homogeneidad evolutiva es la que permite atribuir con confianza los patrones espaciales encontrados (Objetivo Específico 2) a la dinámica del sistema y no a artefactos de la clasificación o del filtrado.

En la fase incipiente, ninguno de los dos campos exhibe todavía una organización definida. El viento se distribuye de forma difusa alrededor del vórtice y la precipitación, aunque ya asimétrica hacia el sector oriental por la influencia temprana del WCB, todavía no diferencia a los sistemas ordinarios de los que devendrán severos. Es en la intensificación cuando ambos campos comienzan a converger sobre el mismo sector activo del ciclón, el flanco cálido-oriental, de modo que, en la Población Global, viento y precipitación alcanzan durante esta fase su momento de mayor co-localización relativa (moda de viento de 12--15~m\,s$^{-1}$ y precipitación de 10--20~mm\,h$^{-1}$, con máximos que se superponen parcialmente en el mismo cuadrante).

Este acoplamiento inicial es exactamente lo que se rompe en los sistemas p90 al llegar a la madurez. El viento se hiperconcentra en el cuadrante noroeste, con densidades de hasta $27\times10^{-4}$ en SBR y modas desplazadas a 20--25~m\,s$^{-1}$ (colas de hasta 30~m\,s$^{-1}$). La precipitación, en cambio, abandona su pico de intensidad temprano (colas de hasta 60~mm\,h$^{-1}$ durante la intensificación) y colapsa hacia magnitudes mínimas en el núcleo (1--4~mm\,h$^{-1}$ en SBR y LPB), mientras sus máximos residuales se reorganizan en la banda arqueada periférica del sureste y sur. La intrusión seca que caracteriza al modelo de Shapiro-Keyser es el mecanismo que explica esta bifurcación. Al erradicar la humedad del núcleo, separa físicamente la fuente del forzante hidrológico de la del forzante cinemático, que permanece anclado al frente en fractura. En el decaimiento, esta arquitectura de cuadrantes opuestos no se disuelve sino que se acentúa, mientras que en la Población Global la señal simplemente se difumina sin reorganizarse geométricamente.

\section{Por qué el promedio no basta: el argumento estadístico}

Ninguna de estas conclusiones habría sido accesible mediante estadística descriptiva simple. Era necesario cuantificar cómo se distribuye la variabilidad espacial y si dicha distribución cambia de naturaleza entre poblaciones (Objetivo Específico 3). El análisis EOF univariado mostró que, en la Población Global, un único modo dominante basta para describir el campo de viento (EOF1 entre 31\% y 57\% de la varianza, acoplado casi linealmente a la vorticidad, con $\rho$ de hasta $-0.89$ en la madurez de SBR). La dinámica del viento ordinario es, en ese sentido, la respuesta pasiva y predecible de un único forzante sinóptico. En los sistemas p90 este único modo deja de ser suficiente, su varianza cae a 25--38\% y la correlación con la vorticidad colapsa ($\rho \approx -0.26$ a $-0.44$), señal inequívoca de que la organización cinemática de un ciclón intenso obedece a procesos de mesoescala que ningún modo dominante único puede capturar por completo. La precipitación exhibe el patrón inverso pero igualmente revelador. Dispersa entre modos en la Población Global (EOF1 apenas 11--16\%), se concentra abruptamente en un solo modo cuando la banda enroscada se consolida en los sistemas p90 en decaimiento (36.5\% en LPB). La severidad, entonces, no dispersa la variabilidad pluviométrica sino que la reorganiza en una geometría más simple y predecible.

El análisis multivariado (MEOF) fue indispensable para dar el paso final, demostrando que estas dos transformaciones, la fragmentación del viento y la concentración de la precipitación, no son fenómenos paralelos sino la misma transición estructural vista desde ambos campos. En la Población Global, el MEOF1 captura 19--27\% de la covarianza conjunta con cargas co-localizadas en el sector oriental, confirmando que viento y precipitación responden al mismo forzante frontal. En el subconjunto p90 esa covarianza conjunta cae a 12--19\% y emerge, en un único modo dominante, un patrón de anticovarianza geométrica. Las cargas de precipitación se concentran al sureste mientras las de viento lo hacen al noroeste. La correlación negativa entre las componentes principales univariadas de ambos campos ($\rho = -0.49$ en LPB durante la madurez) añade la dimensión temporal a esta historia. Los eventos que más proyectan sobre la banda enroscada son, sistemáticamente, los que menos proyectan sobre el núcleo de viento noroeste. Ninguno de los análisis univariados, por separado, habría revelado esta covarianza cruzada. Es la contribución específica del marco multivariado.

\section{Geografía del desacople}

La magnitud de esta reorganización no es uniforme en el espacio. SBR y LPB, sostenidos por flujos de calor latente oceánico y por el aporte de humedad del Chorro de Capas Bajas de Sudamérica (SALLJ), desarrollan las oclusiones más rápidas y, en consecuencia, el desacople geométrico más pronunciado y las bandas enroscadas más intensas. ARG, en cambio, dominada por una dinámica baroclínica clásica y una modulación orográfica andina menos dependiente de procesos diabáticos locales, presenta oclusiones más graduales, un acoplamiento vorticidad-precipitación que persiste por más tiempo, y patrones de separación más difusos aunque topológicamente consistentes con las otras dos regiones. Este gradiente regional, de un desacople abrupto y termodinámicamente forzado en SBR y LPB a uno gradual y dinámicamente controlado en ARG, confirma que la reorganización estructural documentada no es un artefacto estadístico del catálogo combinado, sino la expresión de regímenes físicos genuinamente distintos.

\section{De la estadística a la geometría estructural}

Los prototipos estructurales acoplados (Objetivo Específico 4), que superponen en un mismo panel las PDFe y el EOF1 de ambas variables, son la síntesis visual de todo lo anterior. En la Población Global confirman separaciones modales moderadas (3--5\textdegree) durante las fases tempranas. En el subconjunto p90, en la madurez de SBR, cuantifican una separación entre la moda de precipitación ($\Delta x = 6{,}8^{\circ}$, $\Delta y = -1{,}9^{\circ}$) y la de viento ($\Delta x = -1{,}8^{\circ}$, $\Delta y = 2{,}4^{\circ}$) superior a $8^{\circ}$ euclídeos, que se incrementa a más de $9^{\circ}$ en el decaimiento. Este patrón se repite en LPB y, más atenuado, en ARG. Esta cuantificación geométrica explícita traduce el hallazgo estadístico en un enunciado espacial concreto. En un ciclón intenso maduro del Atlántico Sudoccidental, los máximos de viento y de precipitación no coinciden en el espacio, una asimetría estructural con implicaciones directas para la gestión del riesgo costero.
